\chapter{Inflation, Strukturbildung und der Verlust der Anfangsbedingungen}

\section{Omega und der Nachweis der räumlichen Flachheit}

Die Frage eines Studenten zu großem \(\Omega\) führt zurück zur Friedmann-Robertson-Walker-Gleichung. Für die heutige Epoche, lange nach dem Ende der Strahlungsdominanz, wird das Expansionsgesetz mit einem Term für die Vakuumenergie und einem Materieterm geschrieben. Die Strahlung wird weggelassen, weil ihr Beitrag inzwischen sehr klein ist:
\begin{equation}
H^2=\left(\frac{\dot a}{a}\right)^2=\frac{8\pi G}{3}\left(\rho_0+\rho_m\right)-\frac{k}{a^2}.
\end{equation}
Der Vakuumterm verdünnt sich nicht. Der Materieterm schon:
\begin{equation}
\rho_m=\frac{M}{a^3}.
\end{equation}
\(\Omega\) ist der Anteil der jeweiligen Komponente an der Expansionsgleichung.

Teilt man durch \(H^2\), addieren sich die normierten Terme zu eins. In einer regularisierten Schreibweise gilt
\begin{equation}
\begin{aligned}
\Omega_0&=\frac{8\pi G\rho_0}{3H^2},\\
\Omega_m&=\frac{8\pi G\rho_m}{3H^2},\\
\Omega_k&=-\frac{k}{a^2H^2},
\end{aligned}
\end{equation}
und damit
\begin{equation}
\Omega_0+\Omega_m+\Omega_k=1.
\end{equation}
\editorialnote{Die Tonaufnahme und das Tafelbild sind in diesem Abschnitt schwer verständlich, und der Sprecher verwendet mehrere Bezeichnungen für den Vakuumbeitrag. Die Notation hier regularisiert die Argumentation nach der Division durch \(H^2\).}
Ist der Raum flach, gilt \(k=0\), also
\begin{equation}
\Omega_0+\Omega_m=1.
\end{equation}
Ist die Summe aus Materie und Vakuumenergie größer als eins, ist die Krümmung positiv. Ist sie kleiner als eins, ist die Krümmung negativ.

Die Beobachtungslogik ist dann einfach. Man misst gewöhnliche Materie und Dunkle Materie, die Hubble-Expansion und die Krümmung geometrisch auf sehr großen Skalen. Außerdem stellt man fest, dass das Universum beschleunigt expandiert. Das Vorzeichen des Vakuumterms wird nicht durch \(H\) allein bestimmt, sondern durch die Tendenz von \(\dot a/a\), sich einem konstanten Wert anzunähern. Das bedeutet, dass sich der Skalenfaktor dem exponentiellen Verlauf annähert,
\begin{equation}
a(t)\propto e^{Ht}.
\end{equation}
Die verschiedenen Messungen prüfen einander gegen und zeigen, dass das Universum bis auf wenige Prozent nahezu flach ist.

\begin{classroomqa}
\lecturetimestamp{00:10:50}
\audiencequestion{Umfasst Materie auch Dunkle Materie?}
\lecturerresponse{Ja. Für die gegenwärtige Diskussion umfasst der Materieterm auch Dunkle Materie. Man könnte außerdem Strahlung hinzufügen, aber ihr Anteil ist heute so klein, dass man ihn vernachlässigen kann.}
\end{classroomqa}

\begin{classroomqa}
\lecturetimestamp{00:11:11}
\audiencequestion{Gibt es neben der Hubble-Konstanten noch eine andere Möglichkeit, Dunkle Energie zu messen?}
\lecturerresponse{Entscheidend ist nicht der Wert von \(H\) allein, sondern die Art, wie sich die Expansion einem exponentiellen Gesetz annähert. Zusammen mit der Materiemessung und der geometrischen Messung der räumlichen Flachheit lässt sich so der Vakuumterm erschließen.}
\end{classroomqa}

\section{Von der räumlichen Flachheit zur Strukturbildung}

Warum ist das Universum so flach? Die vorgeschlagene Erklärung ist die Inflation. Flach bedeutet hier: sehr groß, sehr glatt und auf Skalen von Milliarden, ja sogar Dutzenden Milliarden Lichtjahren sehr homogen. Die exponentielle Expansion kann ein Gebiet so weit strecken, dass es diesen Anschein erweckt.

Doch auch Gutes kann man übertreiben. Ein vollkommen glattes Universum bleibt vollkommen glatt. Die Gravitation erzeugt keine Klumpen aus exakter Symmetrie. Wenn es heute Strukturen gibt, müssen sie aus einer kleinen früheren Inhomogenität entstanden sein. Verfolgt man die Entwicklung des heutigen Universums rückwärts, ergibt sich für die Inhomogenität zum Zeitpunkt der Transparenz ungefähr ein Teil in \(10^5\):
\begin{equation}
\frac{\delta\rho}{\rho}\sim 10^{-5}.
\end{equation}
Das Universum war außerordentlich glatt, aber nicht vollkommen.

Von diesem Punkt an bewirkt die Gravitation etwas, was gewöhnliche Kräfte in einem mit Luft gefüllten Raum nicht tun. Eine kleine Überdichte zieht Materie an. In der Umgebung fehlt dann Materie, sodass die Unterdichte noch unterdichter und die Überdichte noch überdichter wird. Die Gravitation verstärkt den Kontrast, statt ihn auszugleichen.

Die naheliegende Frage ist, warum das nicht schon früher geschah. Die Antwort lautet: Druck. Zu früheren Zeiten war der Strahlungsdruck groß. Ausgehend von
\begin{equation}
p=w\rho,
\end{equation}
lauten die bekannten Fälle
\begin{equation}
w_{\rm rad}=\frac13,\qquad
w_m=0,\qquad
w_{\rm vac}=-1.
\end{equation}
Positiver Druck homogenisiert. Er treibt Störungen auseinander und glättet das Medium. Materie, die auf großen Skalen als drucklos behandelt wird, bietet keine solche Stütze. Solange Druck vorhanden ist, wird die Gravitation in Schach gehalten. Wird der Druck vernachlässigbar, setzt sich die Gravitation durch.

Dunkle Materie ist für die Reihenfolge der Ereignisse entscheidend. Die ersten kollabierenden Klumpen sind Halos aus Dunkler Materie. Gewöhnliche baryonische Materie fällt anschließend in diese Halos und sammelt sich, weil Baryonen Energie dissipieren können, näher am Zentrum. Dunkle Materie dissipiert nicht auf dieselbe Weise und bleibt daher im Halo verteilt. Der Beginn dieser Kontraktion liegt etwas vor dem Zeitpunkt der Transparenz, und zwar ungefähr um einen Faktor zehn in der Zeit. Deshalb zeigt der Mikrowellenhintergrund nicht unmittelbar ihren Beginn.

Ein kurzer Exkurs ordnet die spätere chemische Entwicklung ein. Fast alle Elemente, die schwerer als Helium sind, entstehen durch Sternentwicklung und nicht durch die Kosmologie des frühen Universums. Im frühen Universum entstehen im Wesentlichen Wasserstoff, Helium und geringe Spuren einiger leichter Atomkerne.

Beobachtet wird am Himmel nicht die Massendichte selbst, sondern die Temperaturanisotropie des Mikrowellenhintergrunds. Dichtefluktuationen tragen dazu bei, ebenso Geschwindigkeitsfluktuationen. Überdichte Gebiete erscheinen tendenziell etwas kälter, weil Photonen beim Verlassen tieferer Gravitationsmulden zusätzliche Energie verlieren; hinzu kommen Dopplerverschiebungen durch die Bewegung des Gases. Die vollständige numerische Analyse ist kompliziert. Die robuste Aussage ist einfacher: Als das Universum transparent wurde, variierte die Energiedichte um etwa einen Teil in \(10^5\).

\begin{classroomqa}
\lecturetimestamp{00:20:29}
\audiencequestion{Warum ließ die Gravitation die Klumpen nicht schon früher wachsen?}
\lecturerresponse{Weil das Universum zuvor strahlungsdominiert war und die Strahlung genügend positiven Druck aufwies, um das Medium zu homogenisieren. Erst als die Strahlung gegenüber der Materie unbedeutend wurde, begann die Gravitation das Wachstum der Klumpen zu beherrschen.}
\end{classroomqa}

\begin{classroomqa}
\lecturetimestamp{00:24:35}
\audiencequestion{Braucht man für den Kollaps nicht einen Dissipationsmechanismus?}
\lecturerresponse{Gewöhnliche baryonische Materie braucht tatsächlich Dissipation, um sich im Zentrum anzusammeln. Dunkle Materie kann jedoch zunächst Halos bilden, ohne dass eine solche dissipative Wechselwirkung durch Stöße nötig ist. Außerdem gibt es eine kosmologische Dämpfung, die mit der Expansion selbst zusammenhängt und später wichtig wird.}
\end{classroomqa}

\begin{classroomqa}
\lecturetimestamp{00:30:34}
\audiencequestion{Gab es auch Geschwindigkeitsfluktuationen, sodass spätere Galaxien Drehimpuls erwerben konnten?}
\lecturerresponse{Ja. Das primordialen Muster umfasste Dichte- und Geschwindigkeitsfluktuationen. Am Himmel sehen wir ein Temperaturmuster, das sowohl durch Dichteunterschiede als auch durch die Bewegung des Gases entsteht.}
\end{classroomqa}

\section{Das Inflaton und das Plateau}

Der nun eingeführte Mechanismus beruht auf einem Skalarfeld \(\phi\). Seine Energiedichte setzt sich aus drei Beiträgen zusammen:
\begin{equation}
\rho_\phi=\frac12\dot\phi^2+\frac12(\nabla\phi)^2+V(\phi).
\end{equation}
Zeitliche Änderungen tragen ebenso bei wie räumliche Änderungen. Ein Skalarfeld kann außerdem eine echte potenzielle Energie \(V(\phi)\) besitzen. Das ist nichts Exotisches; das Higgs-Feld ist ein bekanntes Beispiel. Entscheidend ist hier die Form des Potenzials.

Das Potenzial muss über einen gewissen Bereich ziemlich flach sein, mit einer leichten Abwärtsneigung; bei einem kritischen Wert fällt der Grund dann steil ab. Er fällt nicht exakt auf null, sondern auf einen sehr kleinen Wert. In Planck-Einheiten liegt die heutige Vakuumenergie in der Größenordnung
\begin{equation}
10^{-123}.
\end{equation}
Auf dem Plateau war sie viel größer, aber immer noch weit unterhalb der Planck-Skala.

\begin{figure}[t]
\centering
\includegraphics[width=0.78\textwidth]{lecture_08_figure_03.png}
\caption{Qualitatives Inflatonpotenzial mit einem breiten Plateau, einem steilen Abfall und einem niedrigen Minimum.}
\lectureframe{00:43:23}
\end{figure}

Diese Darstellung zeigt keinen Ort im Raum, sondern einen Feldwert. Zu sagen, das Universum habe weit oben auf dem Plateau begonnen, bedeutet lediglich, dass der Feldwert in einem sehr großen Gebiet dort lag. Es spielt keine Rolle, wenn \(\phi\) anfangs von Ort zu Ort etwas variierte, denn die Inflation streckt diese Variationen enorm. Nach hinreichender Expansion erscheint jedes lokale Gebiet nahezu gleichförmig. \editorialnote{Die Achsenbeschriftungen sind auf dem Tafelbild nicht lesbar. Aus der mündlichen Erläuterung geht jedoch hervor, dass die vertikale Achse die potenzielle Energie \(V(\phi)\) und die horizontale Achse den Feldwert \(\phi\) darstellt.}

Solange das Feld auf dem breiten, nur leicht geneigten Teil des Potenzials liegt, verhält sich die potenzielle Energie wie Vakuumenergie. Wäre \(\phi\) vollständig eingefroren, würde sie sich exakt wie Vakuumenergie verhalten. Lässt man die nicht angegebenen numerischen Faktoren beiseite, lautet die in der Vorlesung angegebene Beziehung einfach
\begin{equation}
H\propto \sqrt{V(\phi)},
\end{equation}
Daraus folgt das nahezu exponentielle Gesetz
\begin{equation}
a(t)\propto e^{Ht}.
\end{equation}
Fast alles, was heute beobachtet werden kann, geschieht, nachdem das Feld über die Kante gerollt ist. Beim Abfallen wird die große, vakuumartige Energie des Plateaus in Strahlung, gewöhnliche Materie, Dunkle Materie und die geringe Vakuumenergie umgewandelt, die heute noch übrig ist.

\begin{classroomqa}
\lecturetimestamp{00:44:07}
\audiencequestion{Auf welchem Raum ist \(\phi\) ein Feld, und verwendet man weiterhin Einsteins Gleichungen?}
\lecturerresponse{Ja. \(\phi\) ist ein Feld auf der Raumzeit, und dieselbe Friedmann-Gleichung wird weiterhin verwendet. In diesem Bereich wird die Energiedichte auf der rechten Seite hauptsächlich durch \(V(\phi)\) bestimmt, sofern sich das Feld langsam genug ändert.}
\end{classroomqa}

\begin{classroomqa}
\lecturetimestamp{00:47:04}
\audiencequestion{Wie kann sich die vakuumartige Energie des Inflatons in Materie und Strahlung umwandeln?}
\lecturerresponse{Wenn das Feld zum Minimum rollt und oszilliert, entsprechen diese Oszillationen den Quanten des Feldes. Sind die Quanten instabil, können sie in Photonen, Elektronen, Positronen und andere bekannte Teilchen zerfallen.}
\end{classroomqa}

\section{Langsames Rollen und Hubble-Reibung}

Warum sollte das Feld, nachdem es über die Kante gerollt ist, nicht einfach für immer hin und her schwingen? Warum sollte es nicht zurückschwingen und vielleicht sogar wieder hinaufsteigen? Die hier gegebene Antwort ist ein kosmologischer Dämpfungsmechanismus. Es handelt sich nicht um echte mikroskopische Reibung, aber in der Bewegungsgleichung wirkt der Mechanismus auf dieselbe Weise.

Vernachlässigt man räumliche Gradienten und behandelt das Feld als homogen, bilden der kinetische und der potenzielle Term ein System mit einem Freiheitsgrad und \(T-V\)-Dynamik. Da es sich jedoch um Energiedichten handelt, müssen sie mit dem physikalischen Volumen \(a^3\) multipliziert werden. So erhält man eine homogene effektive Lagrange-Funktion der Form
\begin{equation}
L_{\mathrm{hom}}=a^3(t)\left[\frac12\dot\phi^2-V(\phi)\right].
\end{equation}
\editorialnote{Die oberste Zeile an der Tafel ist teilweise abgeschnitten. Diese homogene Lagrange-Funktion wird nur so weit rekonstruiert, wie es nötig ist, um die darunter notierte gesicherte Bewegungsgleichung herzuleiten.}

Die an der Tafel eindeutig erkennbare Gleichung lautet
\begin{equation}
\frac{d}{dt}\left(a^3\dot\phi\right)=-a^3\frac{\partial V}{\partial\phi}.
\end{equation}
Leitet man die Zeitableitung aus, erhält man
\begin{equation}
a^3\ddot\phi+3a^2\dot a\,\dot\phi=-a^3\frac{\partial V}{\partial\phi},
\end{equation}
und durch Division durch \(a^3\) folgt
\begin{equation}
\ddot\phi+3H\dot\phi=-\frac{\partial V}{\partial\phi}.
\end{equation}
Das Minuszeichen wird in der mündlichen Herleitung ausdrücklich korrigiert: Die Kraft ist das Negative der Ableitung der potenziellen Energie.

\begin{figure}[t]
\centering
\includegraphics[width=0.78\textwidth]{lecture_08_figure_05.png}
\caption{Homogene Bewegungsgleichung des Inflatons mit dem Hubble-Reibungsterm \(3H\dot\phi\).}
\lectureframe{00:57:10}
\end{figure}

Der Term \(3H\dot\phi\) wird als Hubble-Reibung bezeichnet. Er entsteht vollständig durch die Ableitung des Faktors \(a^3\). Dynamisch wirkt er jedoch wie der Widerstand, den ein durch Wasser fallender Stein erfährt. Es gibt einen Kraftterm, einen Beschleunigungsterm und einen zur Geschwindigkeit proportionalen Term. Das Feld erreicht rasch eine Endgeschwindigkeit und bewegt sich dann langsam entlang des Plateaus. Das ist das Slow-Roll-Regime. Solange die Steigung gering und die Bewegung des Feldes langsam ist, ändert sich die potenzielle Energie nur langsam, sodass die Expansion nahezu exponentiell bleibt.

In der Nähe des Minimums oszilliert das Feld. Diese Oszillationen entsprechen den Inflatonquanten. Sind die Quanten instabil, zerfallen sie; so wird die Energie des Plateaus in den heißen Materie- und Strahlungsinhalt des späteren Universums umgewandelt.

\begin{classroomqa}
\lecturetimestamp{00:50:48}
\audiencequestion{Warum oszilliert das Feld nicht einfach für immer, und was bewirkt die Dämpfung?}
\lecturerresponse{Die Expansion selbst erzeugt einen reibungsartigen Term. Der Faktor \(a^3\) in der homogenen Lagrange-Funktion führt zum Term \(3H\dot\phi\), der wie viskoser Widerstand wirkt.}
\end{classroomqa}

\section{Quantenfluktuationen und das Überschreiten des Horizonts}

Nun ergibt sich ein scheinbares Problem. Wenn die Inflation das Universum außerordentlich glatt macht, warum löscht sie dann nicht auch die Keime aus, die Galaxien benötigen? Die Antwort lautet: Quantenmechanik. Die überlebenden Fluktuationen sind keine Überreste eines unvollkommenen Glättens. Sie entstehen neu.

Lassen wir die Expansion einen Moment außer Acht. Selbst wenn ein Feld nahezu vollkommen homogen wäre, ließen sich nach der Quantenmechanik \(\phi\) und \(\dot\phi\) nicht beide exakt festlegen. Auch der Grundzustand ist weiterhin in Bewegung. Das ist derselbe Grund, aus dem ein gewöhnlicher harmonischer Oszillator eine Nullpunktsenergie besitzt. Im Alltag sind solche Fluktuationen schwer zu beobachten, weil sie kurzwellige und rasch oszillierende Moden sind, über die Beobachtungen mitteln.

Durch die Expansion ändert sich die Lage. In mitbewegten Koordinaten behält eine Feldmode ungefähr dieselbe Koordinatenwellenlänge, während ihre physikalische Wellenlänge mit dem Skalenfaktor wächst:
\begin{equation}
\lambda_{\rm phys}=a(t)\,\Delta x.
\end{equation}
Während der Inflation bleibt die physikalische Hubble-Skala ungefähr konstant,
\begin{equation}
d_H\sim H^{-1},
\end{equation}
sodass der entsprechende mitbewegte Horizont schrumpft wie
\begin{equation}
d_{H,\mathrm{coord}}\sim \frac{1}{aH}.
\end{equation}

Die Analogie mit einem Seil macht das Einfrieren anschaulich. Eine Welle oszilliert, weil jeder Teil des Seils von benachbarten Teilen gezogen wird. Ist die Wellenlänge jedoch so groß wie der Horizont, liegen die maßgeblichen Nachbarbereiche außerhalb des kausalen Kontakts. Die rücktreibende Kraft verschwindet. Sobald eine Mode länger als der Horizont ist, friert sie ein.

Das geschieht immer wieder. Längere Wellenlängen frieren zuerst ein, kürzere später. Selbst eine Mode, die mit einer winzigen Wellenlänge beginnt, wird schließlich so weit gestreckt, dass sie den Horizont überschreitet. Die Inflation löscht Strukturen also nicht einfach aus, sondern erzeugt fortlaufend eingefrorene Feldvariationen.

Der Vorgang endet auch nicht einfach deshalb, weil die Inflation lange dauert. Ältere eingefrorene Moden werden auf gewaltige Skalen gestreckt, doch ständig frieren neue Moden ein. Das angesammelte Muster wird in dem für die Kosmologie maßgeblichen Sinn statistisch stationär. So entsteht das annähernd skaleninvariante Spektrum: Derselbe Prozess des Horizontübertritts wiederholt sich auf aufeinanderfolgenden logarithmischen Skalen und erzeugt in den verschiedenen Wellenlängenbereichen ungefähr dieselbe Amplitude.

Nachdem das Feld das Plateau verlassen hat, wächst der Horizont wieder. Moden, die zuvor außerhalb des vorübergehenden inflationären Horizonts lagen, treten später wieder ein. Üblicherweise sagt man, dass sie während der Inflation den Horizont verlassen und danach wieder in ihn eintreten. Zu diesem Zeitpunkt sind sie auf makroskopische Skalen gestreckt, sodass sie, wenn sie erneut zu oszillieren beginnen, sehr große Strukturen bilden.

\begin{classroomqa}
\lecturetimestamp{01:12:02}
\audiencequestion{Wie sieht der Horizont im Raumzeitdiagramm aus?}
\lecturerresponse{Während der exponentiellen Expansion bleibt der physikalische Abstand zum Horizont ungefähr bei \(c/H\) konstant. In einem mitbewegten Diagramm erscheint dieser feste physikalische Abstand jedoch immer kleiner, weil die Koordinatenskala gestreckt wird. In mitbewegten Koordinaten schrumpft der Horizont wie \(1/(aH)\).}
\end{classroomqa}

\begin{classroomqa}
\lecturetimestamp{01:15:59}
\audiencequestion{Was geschieht mit der Amplitude, wenn die Fluktuationen gestreckt werden, und gibt es einen stationären Zustand?}
\lecturerresponse{Einzelne Moden werden durch die Streckung verdünnt, doch fortlaufend frieren neue Moden ein. Das Ergebnis ist ein statistisch stationäres, annähernd skaleninvariantes Muster, nicht das vollständige Verschwinden der Fluktuationen.}
\end{classroomqa}

\section{Von eingefrorenen Feldvariationen zu Dichtevariationen}

Sobald das Feld kleine räumliche Variationen aufweist, werden diese in Dichteunterschiede umgewandelt, wenn verschiedene Gebiete zu leicht unterschiedlichen Zeiten über die Kante rollen. Ein Gebiet, das etwas weiter hangabwärts liegt, erreicht die Klippe zuerst. Wird seine Vakuumenergie in Strahlung umgewandelt, verdünnt sich diese Strahlung mit der Expansion. Ein benachbartes Gebiet, das noch auf dem Plateau liegt, bleibt vakuumdominiert und verdünnt sich nicht. Wenn das zweite Gebiet hinunterrollt, ist die Energie des ersten bereits teilweise rotverschoben und damit verringert worden.

Deshalb führt langsameres Rollen zu größeren endgültigen Dichtekontrasten. Bewegt sich das Feld nur sehr langsam entlang des Plateaus, entspricht einem kleinen Unterschied im Feldwert ein größerer Zeitabstand zwischen dem Erreichen der Klippe in verschiedenen Gebieten. Das früher hinunterrollende Gebiet hat mehr Zeit, sich zu verdünnen, bevor das später folgende Gebiet seine Energie umwandelt. Je flacher die Steigung, desto kleiner die Endgeschwindigkeit und desto größer die endgültige Dichtevariation. Ein höheres Plateau führt ebenfalls tendenziell zu größeren Fluktuationen. Die Breite des Plateaus bestimmt hauptsächlich, wie lange die Inflation dauert.

Durch die Fluktuationsdaten ist nur eine Kombination aus Höhe und Steigung eng beschränkt. Ist die Energieskala hoch genug, sollte außerdem ein messbarer Beitrag von Gravitationswellen auftreten. Die numerischen Grenzen werden nur als grobe Richtwerte angegeben. Grob gesagt würde der Mikrowellenhintergrund bereits zu viele Gravitationswellenstrukturen zeigen, wenn die Energiedichte auf dem Plateau deutlich über etwa \(10^{-14}\) in Planck-Einheiten läge. Läge sie deutlich unter etwa \(10^{-16}\), gäbe es kaum eine realistische Aussicht, ein solches Signal nachzuweisen. \editorialnote{Der Sprecher behandelt diese Zahlen ausdrücklich als Näherungswerte und sagt, dass er sich nicht an die genauen Werte erinnert. Sie werden hier nur als Größenordnungsangaben beibehalten.}

Die Einheiten werden wie üblich anhand natürlicher Einheiten erklärt. Die Compton-Wellenlänge eines Teilchens der Masse \(m\) ist
\begin{equation}
\lambda_C=\frac{\hbar}{mc}.
\end{equation}
Wenn \(c=\hbar=1\) gilt, haben Masse und inverse Länge dieselbe Einheit. Die Energiedichte ist Energie pro Volumen und hat daher die Einheit Masse hoch vier:
\begin{equation}
[\rho]=(\mathrm{mass})^4.
\end{equation}

Betrachtet man einen einzelnen Punkt, erscheint die Ansammlung eingefrorener Moden wie ein zufälliger Gang. An verschiedenen Punkten im Raum verlaufen diese zufälligen Gänge unterschiedlich. Genau dadurch wird das Quantenrauschen in eine räumliche Variation umgewandelt.

\begin{classroomqa}
\lecturetimestamp{01:33:04}
\audiencequestion{Führt \(\phi\) einen zufälligen Gang aus, und macht das den früheren Anfangspunkt unwichtig?}
\lecturerresponse{An jedem festen Punkt addieren sich die aufeinanderfolgenden eingefrorenen Moden mit zufälliger Phase, sodass \(\phi\) tatsächlich einen zufälligen Gang ausführt. Gerade dieser Erfolg bringt auch die Schwierigkeit mit sich: Die Inflation ersetzt frühere Informationen fortlaufend durch spätere Fluktuationen, sodass sich die Phase vor der Inflation nur schwer rekonstruieren lässt.}
\end{classroomqa}

\section{Was die Inflation verbirgt}

Die Theorie ist in einem ganz bestimmten Sinn erfolgreich. Sie verknüpft Fluktuationen über einen gewaltigen Skalenbereich durch einen einzigen zugrunde liegenden Mechanismus. Derselbe Mechanismus verwischt jedoch, was zuvor geschah. Frühere Unregelmäßigkeiten werden weggestreckt, und die später erzeugten Quantenfluktuationen sind es, die relevant bleiben. Deshalb ist die Inflation zugleich attraktiv und frustrierend. Sie erklärt die Glattheit und das Fluktuationsspektrum, löscht aber auch einen großen Teil der früheren Geschichte aus.

Das führt zu einem Unbehagen angesichts des Anfangszustands. Eine sehr hohe Energiedichte ist nicht das Problem. Das Problem ist, dass die Energie des Plateaus weder natürlich in der Planck-Größenordnung noch natürlich winzig erscheint. In gewöhnlichen Einheiten ist sie hoch, in Planck-Einheiten aber weiterhin klein. Sie wirkt nicht wie ein naheliegender erster Ausgangspunkt einer fundamentalen Theorie.

In diesem einfachen Bild gibt es auch kein buchstäbliches \(t=0\). Sobald die Inflation zu lange gedauert hat, ist fast jede Information über frühere Phasen verschwunden. Läge die Zahl der E-Faltungen nur knapp über dem erforderlichen Minimum, könnte vielleicht eine Spur der vorherigen Phase erhalten bleiben. Im Allgemeinen löscht die Inflation jedoch die Vergangenheit aus.

Dieselbe Einschränkung gilt für die Behandlung des de-Sitter-Raums. Das rein exponentielle Gesetz beschreibt den inflationären Ast zu späten Zeiten korrekt, doch die exakte de-Sitter-Lösung lässt sich besser durch
\begin{equation}
a(t)\propto \cosh(Ht).
\end{equation}
beschreiben. Diese Lösung kontrahiert bis zu einer minimalen Größe und expandiert anschließend. Ob der kontrahierende Ast in der Kosmologie wörtlich zu verstehen ist, bleibt offen.

\begin{classroomqa}
\lecturetimestamp{01:37:31}
\audiencequestion{Könnte es keinen Anfang gegeben haben, sodass sich das exponentielle Verhalten unbegrenzt weit in die Vergangenheit fortsetzte?}
\lecturerresponse{Nicht in diesem einfachen Bild. Der exakte de-Sitter-Raum beschreibt keine ewig und zu allen Zeiten fortgesetzte exponentielle Expansion, sondern eher \(\cosh(Ht)\) mit einem kontrahierenden und einem expandierenden Ast. Ob der frühere Ast wörtlich zu verstehen ist, bleibt eine offene Frage.}
\end{classroomqa}

\section{Feinabstimmung, Krümmung und anthropische Grenzen}

In erster Näherung sind die wichtigen Parameter die Höhe und die Steigung des Plateaus sowie seine Breite \(\Delta\phi\). Die Steigung ist wichtig, weil sie die Endgeschwindigkeit und damit die Größe der Fluktuationen festlegt. Die Breite ist wichtig, weil sie mitbestimmt, wie lange die Inflation dauert. Eine verbreitete Vorannahme lautet, dass eine grundlegende Feldänderung nicht viel größer als die Planck-Masse sein sollte und dass die Energiedichte die Größenordnung eins in Planck-Einheiten nicht überschreiten sollte. Unter solch groben Grenzen reicht eine ausreichende Inflation nicht für den gesamten Parameterraum. Die hier angegebene Schätzung erfordert eine Feinabstimmung von wenigen Prozent, etwa \(2\%\), damit die Steigung flach genug ist für ungefähr \(60\) bis \(62\) E-Faltungen.

\begin{classroomqa}
\lecturetimestamp{01:46:02}
\audiencequestion{Lässt sich Inflation anhand eines Parameterraums verstehen, und wie fein muss sie abgestimmt sein?}
\lecturerresponse{Ja. In erster Näherung betrachtet man die Höhe, die Steigung und die Breite des Plateaus. Die Vorlesung gibt als grobe Schätzung an, dass eine Feinabstimmung von wenigen Prozent, etwa \(2\%\), nötig ist, um genügend E-Faltungen zu erreichen.}
\end{classroomqa}

Diese mäßige Feinabstimmung wird anschließend mit dem wesentlich schärferen Rätsel der heutigen kosmologischen Konstanten verglichen, deren Größenordnung \(10^{-123}\) beträgt. Die Diskussion wendet sich anthropischen Fragen zu, weil sich beide Probleme neu formulieren lassen, indem man fragt, was für die Entstehung von Strukturen erforderlich ist.

Eine positive kosmologische Konstante wirkt im Newtonschen Bild wie eine universelle, proportional zum Abstand zunehmende Abstoßung. Wäre sie viel größer, müsste sie die Galaxien nicht erst heute auseinanderreißen, um verheerend zu wirken. Es würde genügen, wenn sie groß genug wäre, um zu verhindern, dass sich die primordialen Fluktuationen \(10^{-5}\) jemals zu Galaxien verdichten. Als Beispiel wird angeführt, dass ein Wert, der ungefähr zwei Größenordnungen über dem beobachteten Wert liegt, bereits dazu führen würde. Darauf bezieht sich Weinbergs Argument: Die kosmologische Konstante muss nicht exakt null sein; sie muss lediglich klein genug sein, damit Strukturbildung möglich ist.

Negative Krümmung wirkt in dieselbe Richtung. In der Newtonschen Sprache der früheren Vorlesungen bedeutet \(k<0\), dass die Rezessionsgeschwindigkeit größer als die Fluchtgeschwindigkeit ist. Ein zu starker Auswärtsdrang verhindert ebenfalls, dass sich die kleinen primordialen Überdichten jemals zu Galaxien entwickeln. Die Inflation verdünnt diese Krümmung. Als angegebene untere Grenze gelten ungefähr \(59.5\) E-Faltungen. Ein Wert von etwa \(62\) wird als repräsentative Inflationsgeschichte angenommen, die deutlich über dieser Schwelle liegt.

Dadurch ändert sich die Wahrscheinlichkeitsfrage. Fragt man nach dem reinen Volumen des Parameterraums, das genügend Inflation ermöglicht, beträgt die Antwort möglicherweise nur wenige Prozent. Beschränkt man die Betrachtung auf ein Universum, in dem Galaxien entstehen und Beobachter existieren, die diese Frage stellen können, kann die bedingte Wahrscheinlichkeit wesentlich größer sein. Damit endet die Diskussion: nicht mit einer abschließenden Lösung, sondern mit den Feinabstimmungsrätseln, ihrer anthropischen Neuinterpretation und der Einsicht, dass hier die großen Debatten fortbestehen.
